% Gerado pela etapa `report` (gender_networks.report). Não editar à mão.
\footnotesize
\setlength{\tabcolsep}{4pt}
\begin{tabular}{lrrrrrrrrr}
\toprule
 & \multicolumn{4}{c}{Norma (quantis)} &  & \multicolumn{3}{c}{Dimensões dominantes} &  \\
\cmidrule(lr){2-5}\cmidrule(lr){7-9}
Repr. & 5\% & 50\% & 95\% & máx. & Cos. médio & Maior & $|x|$/mediana & norma² (5 dim.) & $>10\times$med. \\
\midrule
\texttt{lex} & 1,0 & 1,1 & 1,3 & 1,7 & 0,028 & 32 & 2,7 & 0,8\,\% & 0 \\
\texttt{L01} & 6,7 & 8,3 & 11,0 & 21,4 & 0,372 & 0 & 29,6 & 22,5\,\% & 0 \\
\texttt{L18} & 47,3 & 54,8 & 65,0 & 7.636,3 & 0,394 & 4 & 82,0 & 32,7\,\% & 51 \\
\texttt{L36} & 699,2 & 964,1 & 1.198,3 & 1.980,0 & 0,581 & 4 & 33,4 & 25,4\,\% & 0 \\
\texttt{L36n} & 107,9 & 174,7 & 189,8 & 212,7 & 0,593 & 9 & 37,7 & 31,9\,\% & 0 \\
\bottomrule
\end{tabular}

\par\smallskip{\footnotesize\raggedright Cos. médio: cosseno médio entre pares aleatórios de vértices (anisotropia). Maior: índice da dimensão de maior $|x|$ médio; $|x|$/mediana: esse valor sobre a mediana das dimensões; norma²: fração média da norma ao quadrado nas dimensões de maior $|x|$; a última coluna conta vértices com norma acima desse múltiplo da mediana.\par}
